% frontmatter/acknowledgment.tex
\phantomsection % <--- QUAN TRỌNG: Tạo một cái neo ảo tại đây
\begin{center}
\Large \textbf{LỜI CẢM ƠN}
\end{center}

\addcontentsline{toc}{section}{LỜI CẢM ƠN} % Thêm vào mục lục và liên kết với cái neo vừa tạo

    Để hoàn thành bài báo cáo này, em xin gửi lời cảm ơn chân thành đến:

    \begin{itemize}
        \item Ban giám hiệu trường Đại học Tôn Đức Thắng Thành phố Hồ Chí Minh vì đã đưa môn học Quản trị hệ thống mạng vào chương trình giảng dạy, đồng thời tạo điều kiện thuận lợi cho sinh viên về cơ sở vật chất, kiến thức chuyên môn của các giảng viên, giúp cho việc nghiên cứu, học tập thuận lợi hơn.
        \item Đặc biệt, em xin gửi lời cảm ơn giảng viên môn học – thầy Lê Viết Thanh, vì sự giảng dạy tận tâm và chi tiết, truyền đạt những kiến thức quý báu cho em trong suốt quá trình học tập. Nhờ những kiến thức và sự hướng dẫn của thầy, em có thể nắm bắt và hiểu bài học tốt. Đây chắc chắn là những kiến thức quý giá, là nền tảng để em tiếp tục học những môn sau này và áp dụng vào thực tế khi đi làm.
    \end{itemize}

    \begin{itemize}
        \item Môn học Quản trị hệ thống mạng là môn học thú vị, vô cùng hữu ích và có tính tư duy cao. Đảm bảo cung cấp đủ kiến thức, đồng thời còn giúp sinh viên luôn có thể tư duy và tự tìm tòi những phương hướng để ra được kết quả. Do chưa có nhiều kinh nghiệm làm đề tài luận cũng như những hạn chế về kiến thức, trong bài tiểu luận chắc chắn sẽ không tránh khỏi những thiếu sót. Rất mong nhận được sự nhận xét, ý kiến đóng góp từ phía thầy để bài tiểu luận được hoàn thiện hơn.
    \end{itemize}

    Lời cuối cùng, em xin kính chúc thầy nhiều sức khỏe, thành công trên con đường giảng dạy và hạnh phúc trong cuộc sống.

    \vspace{2cm} % Khoảng trống trước chữ ký

\begin{center} % Căn giữa bảng ra giữa trang (hoặc dùng flushright nếu muốn lệch phải)
    \begin{tabular}{c c c} % {c c c} nghĩa là 3 cột đều căn giữa (center)
        \multicolumn{3}{c}{TP. Hồ Chí Minh, ngày 02 tháng 12 năm 2025} \\[0.5cm] % Dòng ngày tháng gộp 3 cột
        
        \textit{Tác giả} & \textit{Tác giả} & \textit{Tác giả} \\
        (Ký tên và ghi rõ họ tên) & (Ký tên và ghi rõ họ tên) & (Ký tên và ghi rõ họ tên) \\[2.5cm] % Khoảng cách để ký tên
        
        \textbf{Nguyễn Tấn Mạnh}& \textbf{Hoàng Hữu Phúc} & \textbf{Phan Trí Tâm}
    \end{tabular}
\end{center}